	\chapter{总结与展望}
\section{总结}
本软件旨在通过自动化和简化的方式提高奖学金评审的效率，减少人为错误，并确保评审过程的公正性和高效性。系统设计涵盖了从启动、界面操作到评审流程的各个环节，具体包括：

\begin{itemize}
	\item 	项目背景：阐述了优秀学生奖学金评审的必要性，并强调了该软件在简化评审流程、提高工作效率方面的重要作用。
	\item 	编写目的：为软件使用者提供详细的操作指南，确保软件能够得到有效使用。
	\item 	适用范围与对象：明确了使用对象为高校评审委员会及相关工作人员，并列举了适用场景，如奖学金评审、学术竞赛评定等。
	\item 	软件概述与界面介绍：展示了软件的启动方式、主界面设计、操作按钮功能等，用户可以通过简单的操作完成学生信息和加分项目的评审。
	\item 	软件操作教程：详细介绍了软件的基本使用流程，涵盖加分项目的数据收集、信息填写、注册与登录操作等。
\end{itemize}
\section{展望}
随着软件版本的更新，以下是未来可能的改进方向：

\begin{enumerate}
	\item 	功能扩展：增加更高级的统计分析功能，如自动生成评审报告、历史数据对比等。
	\item 	界面优化：虽然现有界面已简洁易用，未来可进一步优化界面布局、增加更多自定义选项，提升用户体验。
	\item 	智能化评审：引入人工智能技术，自动识别学生的综合素质评分，减少人工评审的主观性和误差。
	\item 	跨平台支持：目前软件支持Windows系统，未来可以扩展至Mac和Linux平台，增加更多的操作系统兼容性。
\end{enumerate}

\chapter{附录}
考虑到使用说明书的排版美观，本章节用于放置可能出现的长篇幅表格/图像内容（不宜在正文中进行展示或正文中展示排版困难或不美观），可通过点击下行列表中进行快速跳转。

\begin{enumerate}
	\item \hyperref[分值表]{加分项目分值表}
\end{enumerate}


\begin{table}[htbp]
	\centering
	\small
	\renewcommand{\arraystretch}{1} % 调整缩短行间距
	\caption{加分项目分值表}
	\begin{tabular}{cccccc}
		\toprule%第一道横线 
		\textbf{竞赛类加分} & \textbf{加分} & \textbf{科研创新类加分} & \textbf{加分} & \textbf{外语类加分} & \textbf{加分} \\  \midrule%第二道横线            
		国家级一等奖 & 6 & 国家级立项主持人 & 6 & 大学英语四级（CET4） & 1 \\ 
		国家级二等奖 & 4 & 国家级立项学生成员 & 4 & 英语专业四级（TEM4） & 1 \\ 
		国家级三等奖 & 3 & 省级重点立项主持人 & 5 & 雅思（IELTS） & 2 \\ 
		省级一等奖 & 5 & 省级重点立项学生成员 & 3 & 托福（TOEFL） & 2 \\ 
		省级二等奖 & 3 & 省级一般立项主持人 & 3 & 剑桥英语考试 & 2 \\ 
		省级三等奖 & 2 & 省级一般立项学生成员 & 2 & 培生英语考试 & 2 \\
		市级一等奖 & 3 & 校级立项主持人 & 2 & GRE & 2 \\ 
		市级二等奖 & 2 & 校级立项学生成员 & 1 & GMAT & 2 \\ 
		市级三等奖 & 1 & SCI论文一作 & 6 & 其他符合外语类加分项目（2分） & 2 \\
		校级一等奖 & 2 & SCI论文二作 & 4 & 其他符合外语类加分项目（1分） & 1 \\
		校级二等奖 & 1 & SCI论文三作 & 3 & 大学英语六级（CET6） & 2 \\ 
		校级三等奖 & 0.5 & 核心期刊论文一作 & 3 & 英语专业八级（TEM8） & 2 \\ 
		院级一等奖 & 1 & 核心期刊论文二作 & 2 & & \\ 
		院级二等奖 & 0.5 & 核心期刊论文三作 & 1 & & \\ 
		院级三等奖 & 0.3 & 普通期刊论文一作 & 2 & & \\ 
		& & 普通期刊论文二作 & 1 & & \\ 
		& & 普通期刊论文三作 & 0.5 & & \\ 
		& & 版权著作 & 1 & & \\ 
		& & 发明专利一作 & 6 & & \\ 
		& & 发明专利二作 & 4 & & \\ 
		& & 发明专利三作 & 3 & & \\ 
		& & 实用新型专利一作 & 2 & & \\ 
		& & 实用新型专利二作 & 1 & & \\ 
		& & 实用新型专利三作 & 0.5 & & \\ 
		& & 外观专利一作 & 2 & & \\ 
		& & 外观专利二作 & 1 & & \\ 
		& & 外观专利三作 & 0.5 & & \\ 
		\bottomrule%第三道横线
	\end{tabular}
	\label{分值表}
\end{table}